% !TEX program = lualatex
% 日本語版．構成は en/main.tex と同一に保つ．
\documentclass[ja]{lnotes}
% Book-specific preamble, \input by en/main.tex and ja/main.tex (the
% repository root is on TEXINPUTS, so this also works in chapter-only
% builds).  Put this book's own definitions here — listing languages
% (\lstdefinelanguage), shorthand macros, extra packages — instead of
% editing format/lnotes.cls, which is shared by every book.
  % this book's own definitions

\addbibresource{\subfix{../references.bib}}
% 章ごとの文献ファイル（tools/sync-main.py が生成）:
% BEGIN bib (tools/sync-main.py)
% END bib

\title{講義名}
\subtitle{講義ノート}
\author{氏名}
\date{\today}
\subject{COURSE-0000 --- 所属}

\begin{document}
\frontmatter
\maketitle
\tableofcontents

\mainmatter
% 章の一覧．docs/syllabus.json から tools/sync-main.py が生成する．
% 存在する章だけが取り込まれるので，執筆途中でも本全体をビルドできる．
% BEGIN chapters (tools/sync-main.py)
\IfFileExists{lessons/01-using-this-template.tex}{\subfile{lessons/01-using-this-template}}{}
\IfFileExists{lessons/02-lesson-skeleton.tex}{\subfile{lessons/02-lesson-skeleton}}{}
% END chapters

\backmatter
\printbibliography
\printindex
\end{document}
